\chapter{Estado del arte}
\label{chap:estado}

Este capítulo sitúa el trabajo frente a lo publicado sobre tambores de radio variable y mecanismos
de compensación, y frente a las plataformas con las que se ensayan.

\section{Del \emph{fusee} a los tambores de radio variable}

El problema de convertir una fuerza variable en un par constante es antiguo. En los relojes de
muelle, el par del muelle real disminuye a medida que se desenrolla, y eso altera la marcha. El
\emph{fusee} lo resolvió con un cono de perfil curvo y una ranura helicoidal por la que se enrolla
una cadena: con el muelle tenso la cadena tira del radio pequeño, y al debilitarse pasa al radio
grande, de manera que el par entregado al tren se mantiene aproximadamente constante
\cite{landes1983}. La condición de diseño es simplemente que el producto fuerza por radio se
conserve a lo largo del recorrido.

La misma idea se ha formalizado en robótica con el nombre de tambores de radio variable
(\emph{Variable Radius Drums}, VRD). Seriani y Gallina \cite{seriani2016} plantean la síntesis del
perfil de un tambor para que la carga siga una ley prescrita y obtienen soluciones analíticas
cerradas para varios casos. En los robots paralelos accionados por cables, Bieber \emph{et al.}
\cite{bieber2023} diseñan el perfil a partir del perfil de carga máxima del espacio de trabajo, de
modo que el par del motor no supere un valor dado para ninguna longitud de cable. La fabricación
aditiva ha abaratado mucho estos perfiles, lo que explica el interés reciente.

\section{Compensación de gravedad con poleas no circulares}

Un segundo campo de aplicación es la compensación estática de peso. Herder \cite{herder2001} estudia
de forma sistemática los mecanismos estáticamente equilibrados, en los que la energía potencial
total se mantiene constante y el actuador sólo vence las pérdidas. Endo \emph{et al.}
\cite{endo2010} proponen una polea no circular combinada con un muelle para compensar el par
gravitatorio de un brazo, y deducen numéricamente la forma de la polea. Shen y Yano \cite{shen2022}
refinan ese cálculo teniendo en cuenta la variación de la tensión del cable a lo largo del
recorrido, y mejoran así la calidad de la compensación.

El mecanismo de este trabajo comparte el principio, con una diferencia importante: aquí no hay
muelle que almacene energía, sino un motor que aporta todo el trabajo de elevación. La caracola no
ahorra energía, sino que reparte el par.

\section{Bancos de ensayo y validación experimental}

En los trabajos citados, la validación experimental suele consistir en medir el par o la corriente
del actuador a lo largo del recorrido y compararlos con el perfil previsto. Esa comparación exige
tres medidas sincronizadas: posición del eje, fuerza en el cable y par del actuador. Los equipos
comerciales de adquisición que resuelven esto tienen un coste muy superior al del propio prototipo,
y los accionamientos industriales modernos ya publican velocidad y par por bus de campo, de modo que
buena parte de la instrumentación puede sustituirse por software.

Ésa es la oportunidad que aprovecha este TFG: construir, con componentes comerciales de bajo coste y
herramientas libres, un banco que registre las tres magnitudes con la misma base de tiempos, guarde
los datos en bruto y permita contrastar el modelo del mecanismo con la medida. Un resultado
colateral del trabajo, que se detalla en el capítulo~\ref{chap:resultados}, es que esa
instrumentación completa resulta \emph{necesaria}: la medida de fuerza por sí sola no basta para
determinar el perfil del tambor, porque la calibración desconocida de la célula admite soluciones
con el signo invertido; es el par del accionamiento el que deshace la ambigüedad.

\section{Conclusiones del estado del arte}

Los tambores de radio variable están bien estudiados como principio y existen métodos de síntesis de
su perfil. Lo que falta, y es lo que motiva este trabajo, es una plataforma asequible que permita
medir en un prototipo real cómo se reparte el par y contrastarlo con el modelo. Para ello hay que
integrar un accionamiento industrial por Modbus, un acondicionador de célula de carga con reglas de
configuración estrictas y una cadena de adquisición con marcas de tiempo fiables, que es lo que
desarrollan los capítulos siguientes.
